\section{Cálculo Diferencial}

El Cálculo Diferencial surge históricamente de la ambición de Isaac Newton por descifrar las leyes del cosmos. Al estudiar el movimiento de los planetas, Newton enfrentó un desafío que la geometría clásica no podía resolver: las órbitas celestes cambian de dirección y velocidad en cada punto de su trayectoria. Para comprender la mecánica del universo, Newton necesitaba conocer el comportamiento de un cuerpo no en un intervalo de tiempo, sino en un instante preciso.

La solución a este problema fue la concepción de la derivada a través de la \emph{razón de cambio instantánea}. Newton desarrolló un método para capturar la variación de una magnitud en un intervalo tan pequeño que tiende a desaparecer, permitiendo calcular la inclinación exacta de una trayectoria en un \emph{momento infinitesimal}. Geométricamente, este proceso transforma una recta secante en una recta tangente que toca la curva en un único punto, revelando la velocidad exacta del objeto en ese instante.

Hoy estudiamos esta disciplina porque es la herramienta definitiva para la \emph{optimización}. El cálculo nos permite encontrar los puntos máximos y mínimos de cualquier sistema, desde calcular tensiones críticas en ingeniería hasta maximizar ganancias en economía o precisar dosis en medicina. El cálculo diferencial es el lenguaje que Newton utilizó para leer el universo, otorgándonos la capacidad de predecir el comportamiento inmediato de nuestra realidad y entender la \emph{intensidad con la que fluye} todo lo que nos rodea.

\subsection{El problema de la recta tangente}

\subsection{Derivadas de funciones y su cálculo}

\subsection{Derivadas de orden superior}

\subsection{Aplicaciones de la derivada}

\subsubsection{Optimización}
